\section{Convergence analysis}\label{sec:convergence_analysis}

In this section, we study the convergence of Algorithm~\ref{algo:traulls} towards a first-order critical point of problem~\eqref{pb:cnls}. We start by showing global convergence of Algorithm~\ref{algo:trinner_iteration}, implying that the inner minimization phase of Algorithm~\ref{algo:traulls} is always well defined. We then establish global convergence of the main algorithm. Our analysis is based on the criticality measure, i.e. the norm of the projected gradient, which differs from the reduced gradient norm~\eqref{eq:criticality_cond_reduced_grad} used in our implementation. At the end of this section, we discuss the validity of our convergence results when~\eqref{eq:criticality_cond_reduced_grad} is used as a inner iteration termination condition instead of~\eqref{eq:crit_norm_projected_grad}.

\subsection{Global convergence of the inner minimization}

Because our inner-minimization procedure closely resembles the SBMIN algorithm~\cite{conn-etal:1988b} -- used in the bound-constrained inner-minimization phase of LANCELOT~\cite{conn-etal:1992} -- our proof follows the structure of the global convergence proof in~\cite{conn-etal:1988a}. Most of the work lies in reformulating and adapting the intermediate results to the polyhedral case, accounting for the structure of the linear constraints.

We first make the standard assumption
\begin{assumption}\label{assumption:inner_iterates_compact}
	The set $\calX = \left\{x \ | \ \varphi(x) \le \varphi(x_0)\right\} \cap \Omega$ is non empty and compact.
\end{assumption}
By Assumption~\ref{assumption:functions_C2}, it is implicit that the function $\varphi$ is twice continuously differentiable on $\calX$ so we will not state this in the propositions given in this section.

We remind that the quadratic model is of the form
\[q_k(s) = \dfrac{1}{2} \inner{s}{H_ks} + \inner{g_k}{s},\]
where $g_k = \nabla \varphi(x_k)$ and $H_k$ is an approximation of the Hessian based on the structured SR1 formula from~\ref{subsec:hessian_approx}. We formulate two assumptions relative to those approximations. The norm used is the induced norm on matrices, i.e. $\|M\| := \sup_{\|x\|=1} \|Mx\|$ for a given matrix $M$.
\begin{assumption}\label{assumption:model_hessian}
	Defining, for every iteration index $k$, the scalars $b_k$ by 
	\[b_k = 1+\max_{0 \le i \le k} \ \left\| H_i\right\|,\]
	we require that the series $\sum_k \frac{1}{b_k}$ diverges to $\infty$.
\end{assumption}
The second assumption states that the norm of the approximating Hessians should not increase too fast compared with the speed of convergence of the function values.
\begin{assumption}\label{assumption:hessian_norm_compared_convergence_speed}
	\[\lim\limits_{k\to \infty} b_k (\varphi_{k+1}-\varphi_k) = 0.\]
\end{assumption}

The projected gradient path is defined as
\[s_k(t)=\proj{\Omega}{x_k-tg_k}-x_k \text{ for } t\ge 0.\]
The reduction of the model along the projected gradient path may thus be defined as the piecewise quadratic function
\[\psi(t) = q_k(s_k(t)),\]
and we denote by $t_k^C$ the first local minimum of $\psi$ subject to the trust region constraint 
\begin{equation*}
	\|s_k(t)\|_\infty \le \Delta_k.
\end{equation*}
The associated Cauchy step is 
\begin{equation*}
	s_k^C = s_k(t_k^C).
\end{equation*}
Because of the choice of the $\ell_\infty$ norm, computing $s_k(t)$ within the trust region corresponds to project the direction $x_k-tg_k$ onto
\begin{equation}\label{eq:proj_grad_feasible_set}
	\left\{d \in \RR^n \ | \ Ad=0,\ l_k \le d \le u_k\right\},
\end{equation}
with $(l_k)_i = \max\left(-\Delta_k,l_i-(x_k)_i\right)$ and $(u_k)_i = \min\left(\Delta_k,u_i-(x_k)_i\right)$ for all $i=1,\ldots,n$.

We assume that the total step $s_k$ produces a fraction of the reduction achieved by the Cauchy point:
\begin{equation}\label{eq:step_drecrease_wrt_cauchy}
	q_k(s_k) \le \kappa_{fcd} \ q_k(s_k^C),
\end{equation}
for $\kappa_{fcd} \in (0,1]$. 

We detail the behavior of the polygonal line $s_k(t)$. Let $I(t)$ be the set of bounds $l_k$ or $u_k$ satisfied with equality at $s_k(t)$. Notice that this definition slightly differs from the active set $I(x)$ defined earlier, as the new formulation know takes into account the trust region. At first, if no bounds are active at $x_k$, projecting on the set~\eqref{eq:proj_grad_feasible_set} for $t$ close to $0$ is equivalent to projecting onto the null space of $A$. As $t$ increases, the direction might hit several bounds. Since the set of active bounds can only be increased, we have that
\begin{equation}\label{eq:nested_active_sets}
	I(t) \subseteq I(t^\prime) \quad \text{for all}\ 0 < t \le t^\prime.
\end{equation}
Let \(0=t_0 < t_1 < \ldots < t_p\) be the successive values of $t$ at which the projected step $s_k(t)$ hits a bound, also called breakpoints. Hitting a bound not only affects the corresponding components, that are now fixed, but it also affects the other components, as the steepest direction must now be projected on a subspace
\[\left\{d \in \RR^n \ | \ Ad=0,\ d_i=0\ \text{fixed components}\ i\right\}\]
This motivates to adapt the notion of tangent space at a point on the projected gradient path:
\begin{equation*}
	T(t) := \left\{ d \in \RR^n \ | \ Ad = 0,\ d_i=0\ \text{for all}\ i\in I(t)\right\},
\end{equation*}
for $t\ge0$. This enables us to give a recursive expression of $s_k(t)$ on each interval $[t_i,t_{i+1})$, with $0 \leq i \le p$, as:
\begin{equation}\label{eq:projected_gradient_path_recursion}
	s_k(t) = (t-t_i) \proj{T(t_i)}{-g_k} + s_k(t_i).
\end{equation}
We also define the reduced gradient on the projected gradient path:
\begin{equation}\label{eq:reduced_gradient_proj_grad_path}
	z_k(t) := \proj{T(t)}{g_k}.
\end{equation}
Note that the reduced gradient, and hence the path $s_k(t)$ are well defined because of the full rank Assumption~\ref{assumption:full_rank_A} of the matrix $A$. As for the differentiability of $\varphi$, we will not remind it in the results of this section.

We now state a first lemma on how the tangent spaces and the reduced gradients at two different positions compare to each other.
\begin{lemma}\label{lemma:non_increasing_reduced_gradient_norm}
	For all $0 < t < t^\prime$, we have that 
	\begin{equation*}
		T(t) \supseteq T(t^\prime),
	\end{equation*}
	and
	\begin{equation}\label{eq:lemma_reduced_gradient_norm}
		\|z_k(t)\| \ge \|z_k(t^\prime)\|.
	\end{equation}
\end{lemma}
\begin{proof}
	The first statement follows from the inclusion~\eqref{eq:nested_active_sets}. Inequality~\eqref{eq:lemma_reduced_gradient_norm} then results from the fact that projecting the same vector on a smaller linear subspace reduces the norm of its projection. 
\end{proof}
When referring to the tangent space and the reduced gradient at a given breakpoint $t_i$, we will make use of the shorthand notation $T_i := T(t_i)$ and $z_i:=z_k(t_i)$. Notice that because the active set is fixed on each interval $[t_i,t_{i+1})$, so is the tangent space and hence, the reduced gradient is constant. Also, we can deduce from Lemma~\ref{lemma:non_increasing_reduced_gradient_norm} that the tangent spaces associated to each breakpoint form a finite sequence of nested linear subspaces
\[T_0 \supseteq T_1 \supseteq \ldots \supseteq T_p.\]
We can now expand equation~\eqref{eq:projected_gradient_path_recursion}:
\begin{equation}\label{eq:projected_gradient_path_full_expr}
	s_k(t) = -(t-t_i)z_i - \sum_{j=0}^{i-1}(t_{j+1}-t_j)z_j,
\end{equation}
which shows that on each interval $(t_i,t_{i+1})$, $s_k(t)$ is differentiable w.r.t. $t$ and that
\begin{equation}\label{eq:projected_gradient_path_derivative}
	s_k^\prime(t) = -z_k(t) = -z_i,
\end{equation}
for $t \in (t_i,t_{i+1})$.

We start by establishing an inequality on the decrease of the objective function after taking step $s_k$. This will involve the norm of the projected gradient, i.e.:
\begin{equation}\label{eq:crit_norm_projected_grad}
	h_k := \left\|s_k(1)\right\|.
\end{equation}

\begin{lemma}\label{lemma:majoration_reduced_gradient_norm}
	If Assumption~\ref{assumption:inner_iterates_compact} holds and that $h_k>0$, then 
	\[\|z_k(t_k^{(1)})\| \ge \dfrac{h_k}{2},\]
	where \[t_k^{(1)}=\frac{h_k}{2\kappa_{ubg}},\]
	and $\kappa_{ubg}$ is the constant defined by $\kappa_{ubg} := \max\left(1, \max_{x \in \calX} \|\nabla \varphi(x)\|\right)$.
\end{lemma}

\begin{proof} 
	First note that the constant $\kappa_{ubg}$ is well defined because $\varphi$ is continuously differentiable on $\calX$, the latter being compact by Assumption~\ref{assumption:inner_iterates_compact}.
	
	Since the set $\Omega$ is closed and convex, the projection mapping $\proj{\Omega}{\cdot}$ is non expansive, thus, for any $t\ge 0$:
	\begin{equation*}
		\begin{split}
			\left\| s_k\left(t\right) \right\| & = \left\| \proj{\Omega}{x_k-tg_k}-x_k\right\| \\
			& = \left\| \proj{\Omega}{x_k-tg_k}-\proj{\Omega}{x_k}\right\| \\
			& \le \left\| x_k-tg_k-x_k\right\|\\
			& \le t\|g_k\|.
		\end{split}
	\end{equation*}
	Choosing $t=t_k^{(1)}$ and because $\|g_k\|\le \kappa_{ubg}$ by definition of $\kappa_{ubg}$, we get
	\begin{equation}\label{eq:norm_proj_grad_t1}
		\| s_k(t_k^{(1)})\| \le \dfrac{h_k}{2}.
	\end{equation}
	Now, defined $t_k^{(2)}$ as the smallest $t\ge 0$ such that
	\begin{equation}\label{eq:def_t2}
		\|s_k(t_k^{(2)})\| = h_k.
	\end{equation}
	By~\eqref{eq:norm_proj_grad_t1}, we have
	\begin{equation}\label{eq:t1_leq_t2}
		0 < t_k^{(1)} < t_k^{(2)} \le 1.
	\end{equation}
	On each interval $(t_i,t_{i+1})$, the left, resp. right, derivative of $s_k(t)$ at $t_i$, resp. $t_{i+1}$, are well defined and equal $z_i$. We can thus rewrite~\eqref{eq:projected_gradient_path_full_expr} as
	\begin{equation*}
		s_k(t) = \int_{t_i}^{t} s_k^\prime(t)dt + \sum_{j=0}^{i-1} \int_{t_j}^{t_{j+1}} s_k^\prime(t)dt,
	\end{equation*} 
	which we simplify by
	\begin{equation}\label{eq:projected_gradient_path_integral_form}
		s_k(t) = -\int_{0}^{t} z_k(t)dt,
	\end{equation}
	using~\eqref{eq:projected_gradient_path_derivative} and keeping in mind that it is a sum of integrals defined on each segment $[t_i,t_{i+1}]$. 
	Now let $i_1$, resp. $i_2$ be the breakpoint index such that $t_k^{(1)} \in [t_{i_1},t_{i_1+1})$, resp. $t_k^{(2)} \in [t_{i_2},t_{i_2+1})$. Then by~\eqref{eq:projected_gradient_path_integral_form}:
	\begin{equation*}
		s_k(t_k^{(2)}) - s_k(t_k^{(1)}) = -\int_{t_k^{(1)}}^{t_k^{(2)}} z_k(t)dt, 
	\end{equation*}
	which leads to
	\begin{equation*}
		\begin{split}
			\|s_k(t_k^{(2)}) - s_k(t_k^{(1)})\| & \le \int_{t_k^{(1)}}^{t_k^{(2)}} \|z_k(t)\|dt \\
			& \le \int_{t_k^{(1)}}^{t_k^{(2)}} \|z_k(t_k^{(1)})\|dt \\
			& \le (t_k^{(2)}-t_k^{(1)}) \|z_k(t_k^{(1)})\|,
		\end{split}
	\end{equation*}
	where we have used~\eqref{eq:lemma_reduced_gradient_norm} to bound the norm of the reduced gradient on $[t_k^{(1)},t_k^{(2)}]$.
	Combined with~\eqref{eq:t1_leq_t2} and~\eqref{eq:def_t2}, we get
	\begin{equation*}
		\begin{split}
			\dfrac{h_k}{2} = \|s_k(t_k^{(2)})\| - \dfrac{h_k}{2} &\le  \|s_k(t_k^{(2)})\| - \|s_k(t_k^{(1)})\| \\ 
			& \le \|s_k(t_k^{(2)}) - s_k(t_k^{(1)})\| \\
			& \le 	(t_k^{(2)}-t_k^{(1)}) \|z_k(t_k^{(1)})\| \\
			& \le  \|z_k(t_k^{(1)})\|,
		\end{split}
	\end{equation*}
	which is the desired inequality.
\end{proof}

The next lemma gives an upper bound on the quadratic model $\psi(t)$ in an interval of interest.

\begin{lemma}\label{lemma:bound_model_proj_grad_path}
	If Assumption~\ref{assumption:inner_iterates_compact} holds and that for some $t_k^{(3)}>0 $, one has
	\[\alpha_k = \|z_k(t_k^{(3)})\| > 0,\]
	then, if $T$ is the set of points in $[0,t_k^{(3)}]$ at which the piecewise quadratic $\psi$ is differentiable, 
	\begin{equation}\label{eq:lemma_ineq_model_derivative_t3}
		\psi^\prime(t) \le -\alpha_k^2 + t_k^{(3)}\kappa_{ubg}^2\|H_k\|,
	\end{equation}
	for all $t\in T$.
	
	Furthermore, 
	\begin{equation*}
		\psi^\prime(t) \le -\dfrac{1}{2}\alpha_k^2 ,
	\end{equation*}
	for all $t\in T \cap [0,t_k^{(4)}]$ and 
	\begin{equation*}
		\psi(t) \le  - \dfrac{\alpha_k^2}{2}t,
	\end{equation*}
	for all $t\in  [0,t_k^{(4)}]$ where 
	\begin{equation}\label{eq:t4_def}
		t_k^{(4)} = \min \left(t_k^{(3)}, \dfrac{\alpha_k^2}{2\kappa_{ubg}^2(\|H_k\|+1)}\right).
	\end{equation}
\end{lemma}

\begin{proof}
	For $t \in T$, let $i$ be the breakpoint index such that $t\in(t_i,t_{i+1})$. On the latter interval, $\psi$ is differentiable and we have, by expression~\eqref{eq:projected_gradient_path_full_expr}:
	\begin{equation*}
		\begin{split}
			\psi^\prime(t) & = \inner{g_k}{s_k^\prime(t)} + \inner{s_k(t)}{H_ks_k^\prime(t)} \\
			& = -\inner{g_k}{z_i} - \inner{s_k(t)}{H_kz_i}.
		\end{split}
	\end{equation*}
	Because $z_i$ is, by definition, the orthogonal projection of the vector $g_k$ on a linear subspace:
	\begin{equation}\label{eq:ineq_derivative_model_first_order}
		\begin{split}
			\inner{g_k}{z_i} & = \inner{z_i}{z_i} \\
			& \ge \|z_k(t_k^{(3)})\|^2 = \alpha_k^2,
		\end{split}
	\end{equation}
	where the last inequality follows from the monotonicity of the reduced gradient norm on $[0,\infty)$. 
	
	We now look at the quadratic terms. First, by Cauchy-Schwarz inequality:
	\begin{equation*}
		\left|\inner{s_k(t)}{H_kz_i}\right| \le \|s_k(t)\| \|H_kz_i\|.
	\end{equation*}
	Then, applying the triangle inequality to expression~\eqref{eq:projected_gradient_path_full_expr}, we obtain
	\begin{equation}\label{eq:norm_proj_grad_step_ineq}
		\begin{split}
			\|s_k(t)\| & \le (t-t_i) \|z_i\| + \sum_{j=0}^{i-1} (t_{j+1}-t_j)\|z_j\| \\
			& \le (t-t_i) \|g_k\| + \sum_{j=0}^{i-1} (t_{j+1}-t_j)\|g_k\| \\
			& \le t\|g_k\| \le t_k^{(3)}\kappa_{ubg},
		\end{split}
	\end{equation}
	where we have used the facts that for all breakpoints indices $j$, $\|z_i\| \le \|g_k\| \le \kappa_{ubg}$ and that the scalar $t$ is taken in $T$. For the remaining terms:
	\begin{equation*}
		\|H_kz_i\| \le \|H_k\| \|z_i\| \le \kappa_{ubg}\|H_k\|.
	\end{equation*}
	Combining the latter inequality with~\eqref{eq:norm_proj_grad_step_ineq}, we get the following bound for the quadratic terms:
	\begin{equation}\label{eq:ineq_derivative_model_second_order}
		\left|\inner{s_k(t)}{H_kz_i}\right| \le t_k^{(3)} \kappa_{ubg}^2 \|H_k\|.
	\end{equation}
	Hence, using~\eqref{eq:ineq_derivative_model_first_order} and~\eqref{eq:ineq_derivative_model_second_order}:
	\begin{equation*}
		\psi^\prime(t) \le -\alpha_k^2 + t_k^{(3)} \kappa_{ubg}^2 \|H_k\|,
	\end{equation*}
	for all $t \in T$, which proves~\eqref{eq:lemma_ineq_model_derivative_t3}.
	Now, considering $t_k^{(4)}$ as defined by~\eqref{eq:t4_def}, since $t_k^{(4)} \le t_k^{(3)}$, we get that $\|z_k(t_k^{(4)})\| \ge \alpha_k > 0$. The above reasoning based on $t_k^{(4)}$ thus yields
	\begin{equation*}
		\psi^\prime(t) \le -\alpha_k^2 + t_k^{(4)} \kappa_{ubg}^2 \|H_k\| \le -\dfrac{\alpha_k^2}{2},
	\end{equation*}
	for $t\in T$, $t\le t_k^{(4)}$. It follows that, for all $t\in [0,t_k^{(4)}]$:
	\begin{equation*}
		\psi(t) \le -\dfrac{\alpha_k^2}{2}t,
	\end{equation*}
	which completes the proof.
\end{proof}

We can now establish a bound on the decrease guaranteed by the step.
\begin{theorem}
	If assumptions~\ref{assumption:inner_iterates_compact}, \ref{assumption:model_hessian} hold and that $h_k > 0$, then
	\begin{equation}\label{eq:bound_step_decrease}
		q_k(s_k) \le -\kappa_{mdc} h_k^2 \min\left(\dfrac{h_k^2}{b_k},\Delta_k\right),
	\end{equation}
	where 
	\begin{equation*}
		\kappa_{mdc} = \dfrac{\kappa_{fcd}}{64\kappa_{ubg}^2}.
	\end{equation*}
	Furthermore, if iteration $k$ is successful, then
	\begin{equation}\label{eq:bound_successful_step_decrease}
		\varphi(x_k) - \varphi(x_k+s_k) \ge \kappa_{sdc} h_k^2 \min\left(\dfrac{h_k^2}{b_k},\Delta_k\right),
	\end{equation}
	with $\kappa_{sdc} = \eta_1 \kappa_{mdc}$.
\end{theorem} 
\begin{proof}
	We first observe that if we use $t_k^{(1)}$ as $t_k^{(3)}$ in the proof of Lemma~\ref{lemma:bound_model_proj_grad_path} and apply Lemma~\ref{lemma:majoration_reduced_gradient_norm}, we get 
	\begin{equation}\label{eq:model_bound_t1}
		\psi(t) \le -\dfrac{h_k^2}{8} t,
	\end{equation}
	for $t \in [0,t_k^{(5)}]$ with \[t_k^{(5)} := \min\left(t_k^{(1)},\dfrac{h_k^2}{8\kappa_{ubg}^2 (\|H_k\|+1)}\right) = \dfrac{h_k^2}{8\kappa_{ubg}^2 (\|H_k\|+1)}.\]
	We will bound the decrease obtained with the step depending on where the point associated with $t_k^{(5)}$ lies with respect to the trust region boundary.
	
	First, assume that $\|s_k(t_k^{(5)})\|_\infty \le \Delta_k$. By~\eqref{eq:model_bound_t1}, one has $\psi^\prime(t) < -h_k^2 /8 < 0$ for all $t \in  [0,t_k^{(5)}]$, so $\psi$ in strictly decreasing on this segment. Because the Cauchy point is the first local minimizer of $\psi$ subject to the trust region and the trust region is not violated at $t_k^{(5)}$, we must have $t_k^C \ge t_k^{(5)}$. Therefore, 
	\[\psi(t_k^C) \le \psi(t_k^{(5)}) \le -\dfrac{h_k^2}{8} t_k^{(5)} \]
	
	Then, by~\eqref{eq:step_drecrease_wrt_cauchy}, $q_k(s_k) \le \kappa_{fcd} q_k(s_k^C)$ and it follows that
	\begin{equation}\label{eq:model_decrease_cauchy_in_tr}
		q_k(s_k) \le -\kappa_{fcd} \frac{h_k^2}{8} t_k^{(5)}.
	\end{equation}
	Now, assume that $\|s_k(t_k^{(5)})\|_\infty > \Delta_k$. The latter implies that $\|s_k(t_k^C)\|_\infty = \Delta_k$ and from
	\begin{equation*}
		\left\|s_k\left(\frac{\Delta_k}{\kappa_{ubg}}\right)\right\|_\infty \le 	\left\|s_k\left(\frac{\Delta_k}{\kappa_{ubg}}\right)\right\| \le \frac{\Delta_k}{\kappa_{ubg}} \|g_k\| \le \Delta_k,
	\end{equation*}
	we can deduce that 
	\begin{equation*}
		t_k^C \ge \frac{\Delta_k}{\kappa_{ubg}}.
	\end{equation*}
	Therefore, using~\eqref{eq:model_bound_t1} with $t=t_k^C$ and~\eqref{eq:step_drecrease_wrt_cauchy} implies
	\begin{equation}\label{eq:model_decrease_cauchy_at_tr}
		q_k(s_k) \le -\kappa_{fcd} \frac{h_k^2}{8\kappa_{ubg}} \Delta_k \le -\kappa_{fcd} \frac{h_k^2}{64\kappa_{ubg}^2} \Delta_k,
	\end{equation} 
	because $\kappa_{ubg} \ge 1.$
	The inequality~\eqref{eq:bound_step_decrease} results from gathering~\eqref{eq:model_decrease_cauchy_in_tr},~\eqref{eq:model_decrease_cauchy_at_tr} and the definition of scalar $b_k$. 
	
	Finally, if the step is successful, we get inequality~\eqref{eq:bound_successful_step_decrease} by using~\eqref{eq:bound_step_decrease} and the step acceptance condition $\rho_k > \eta_1$ from Algorithm~\ref{algo:trinner_iteration}. 
\end{proof}

Once we have established the guaranteed decrease inequality~\eqref{eq:bound_step_decrease}, the rest of the proof does not involve the polyhedral structure of the constraints and we can fall back into the standard convergence theory of trust region methods. We state the main convergence theorem, whose proof and intermediate developments can be found in~\cite{conn-etal:1988a}.
\begin{theorem}[Theorem 11 from~\cite{conn-etal:1988a}]
	If assumptions~\ref{assumption:inner_iterates_compact}, \ref{assumption:model_hessian}, \ref{assumption:hessian_norm_compared_convergence_speed} hold, then
	\[\lim\limits_{k\to \infty} h_k = 0.\]
\end{theorem}

\subsection{Global convergence of the algorithm}

The content of this section follows the structure of the global convergence proof outlined in~\cite{conn-etal:1996b} for AL algorithms designed to solve problems whose constraints combine nonlinear equalities and linear inequalities. We first make the following assumption about the iterates considered.

\begin{assumption}\label{assumption:iterates_domain_bounded}
	The iterates produced by Algorithm~\ref{algo:traulls} lie within a closed bounded domain of $\Real^n$.
\end{assumption}

Let $(x_k)_k$ be a sequence of such iterates, then Assumption~\ref{assumption:iterates_domain_bounded} ensures the existence of a subsequence indexed by $\calK \subseteq \NN$ converging to a limit point $x_*$. Since the iterates satisfy $x_k \in \Omega$ for all $k$, we also have $x_* \in \Omega$. We recall from~\eqref{eq:tangent_space} that $I(x)$ is the set of bounds active at $x$ and $T(x)$ the associated tangent space, and we abbreviate $I_* := I(x_*)$. We denote by $\tilde A_*$ the matrix gathering the linear constraints active at $x_*$,
\begin{equation*}
	\tilde A_* := \begin{pmatrix} A \\ Z_* \end{pmatrix},
\end{equation*}
where $Z_*$ is the submatrix of the identity whose rows are $e_i^\top$, $i \in I_*$. Consistently with the well-posedness of the projector~\eqref{eq:minor_subpb_projector}, $\tilde A_*$ has full row rank, so $T(x_*) = \mathrm{null}(\tilde A_*)$ and the columns of the orthonormal matrix $N_*$ form a basis of $T(x_*)$. The next assumption ensures that the null space of the active linear constraints is large enough to achieve feasibility of the nonlinear constraints, and that the gradients of those constraints are linearly independent within that space.

\begin{assumption}\label{assumption:cq_null_space_jacobian}
	The rank of the matrix $C(x_*)N_*$ is no smaller than $n_c$ at any limit point $x_*$ of the sequence $(x_k)_k$.
\end{assumption}

We now introduce the remaining notations and notions used in our convergence analysis.

Since $\Omega$ is of the form~\eqref{eq:linear_constraints}, the normal cone to $\Omega$ at $x \in \Omega$ can be defined as
\begin{equation}\label{eq:normal_cone}
	\calN_\Omega(x) = \bigg\{ A^\top\xi + \sum_{i \in I^+(x)} \sigma^+_i e_i - \sum_{i \in I^-(x)} \sigma^-_i e_i \ \Big|\ \xi\in\RR^m,\ \sigma^+_i, \sigma^-_i\ge0\bigg\},
\end{equation}
where the equality multipliers $\xi$ are free in sign while the bound multipliers are signed. The first-order criticality condition $\proj{\Omega}{x_*-\nabla_x\calL(x_*,\lambda_*)}=x_*$ of problem~\eqref{pb:cnls} is then equivalent to
\[-\nabla_x\calL(x_*,\lambda_*)\in \calN_\Omega(x_*).\]

Finally, for vectors $x$ at which the matrix $C(x)N_* N_*^\top C(x)^\top$ is non-singular, the least-squares multipliers associated with $\tilde A_*$ are defined by
\begin{equation*}
	\lambda(x) = -\big((C(x)N_*)^\dagger\big)^\top N_*^\top\nabla f(x),
\end{equation*}
where $(C(x)N_*)^\dagger = N_*^\top C(x)^\top\big(C(x)N_* N_*^\top C(x)^\top\big)^{-1}$ is the right pseudoinverse of $C(x)N_*$. When $(C(x)N_*)^\dagger $ is well defined, Assumption~\ref{assumption:functions_C2} ensures that $\lambda(x)$ is differentiable and that its Jacobian $\nabla\lambda(x)$ is bounded in a neighborhood of $x_*$ (see~\cite[Lemma 2.1]{conn-etal:1996b}).

We can now proceed with the global convergence proof, starting with a first lemma giving a bound on the reduced gradient norm.
\begin{lemma}\label{lemma:bound_proj_algrad_sol_nullspace}
	Let $\left(x_k\right)_{k\in \calK} \in \Omega$ be a sequence converging to $x_*$ and such that
	\[\left\| \proj{\Omega}{x_k - \nabla_x \Phi_k}-x_k \right\| \le \omega_k,\]
	with $\lim\limits_{k\to \infty} \omega_k = 0$. Then
	\begin{equation*}
		\left\| \proj{T(x_*)}{\nabla_x \Phi_k} \right\| \le  \omega_k,
	\end{equation*}
	for all $k\in \calK$ large enough.
\end{lemma}

\begin{proof}
	To lighten the notations, we write $v_k = \proj{\Omega}{x_k - \nabla_x \Phi_k}$ and $d_k = v_k - x_k$. We have
	\begin{equation*}
		\begin{split}
			\|v_k - x_*\| &\le \|v_k - x_k\| + \| x_k - x_*\| \\
			& \le \omega_k + \| x_k - x_*\|.
		\end{split}
	\end{equation*}
	Since, by hypothesis, $\omega_k \to 0$ and $x_k \to x_*$ for $k\in \calK$, the two terms on the right hand side converge to $0$, which implies that we also have
	\begin{equation*}
		\lim\limits_{k \in \calK} v_k = x_*.
	\end{equation*}
	Therefore, there is a threshold index $k_0$ such that all bounds inactive at $x_*$ are also inactive at $v_k$ for $k \ge k_0$. This is equivalent to $I(v_k) \subseteq I_*$ for $k\ge k_0$. We then have the reverse inclusion for the associated tangent spaces $T(x_*) \subseteq T(v_k)$, giving
	\begin{equation}\label{eq:lemma_bound_redgrad_comp_norm_tanproj}
		\left\| \proj{T(x_*)}{\nabla_x \Phi_k} \right\| \le \left\| \proj{T(v_k)}{\nabla_x \Phi_k} \right\|,
	\end{equation}
	whenever $k\in \calK$ is above the threshold index $k_0$.
	
	But recall from Section~\ref{sec:algorithm} that $T(v_k)$ is the null space of a matrix $\tilde A_k$ of the form~\eqref{eq:matrix_tangent_space}. Moreover, since $v_k$ is a projection onto a convex set, it can be characterized by
	\begin{equation*}
	x_k - \nabla_x \Phi_k - v_k = -\left(\nabla_x \Phi_k + d_k\right) \in \calN_\Omega(v_k).
	\end{equation*} 
	By the expression~\eqref{eq:normal_cone}, $\calN_\Omega(v_k)$ is trivially a subset of the row space of $\tilde A_k$ and the latter is also the orthogonal complement of the null space of $\tilde A_k$, i.e.  $T(v_k)^\perp$. 
	Because $T(v_k)^\perp$ is a vector space:
	\[\nabla_x \Phi_k + d_k \in T(v_k)^\perp,\]
	and we obtain, by linearity of $\proj{T(v_k)}{\cdot}$:
	\begin{equation}\label{eq:lemma_bound_redgrad_proj_equality}
		\proj{T(v_k)}{\nabla_x \Phi_k + d_k} = 0 \iff \proj{T(v_k)}{\nabla_x \Phi_k} = -\proj{T(v_k)}{d_k}.
	\end{equation}
	Combining~\eqref{eq:lemma_bound_redgrad_comp_norm_tanproj} and~\eqref{eq:lemma_bound_redgrad_proj_equality} gives
	\begin{equation*}
		\left\| \proj{T(x_*)}{\nabla_x \Phi_k} \right\| \le \left\| \proj{T(v_k)}{\nabla_x \Phi_k} \right\| = \left\| \proj{T(v_k)}{d_k} \right\| \le \|d_k\| \le \omega_k,
	\end{equation*}
	whenever $k \in \calK$ is above the threshold $k_0$, which is the desired inequality.
\end{proof}

The next lemma is a standard result about the behavior of the Lagrange multipliers in AL methods.
\begin{lemma}[Lemma 14.4.1~\cite{conn-etal:2000}]\label{lemma:behavior_multipliers}
	Assume that $\lim\limits_{k\to \infty} \mu_k = \infty$ when Algorithm~\ref{algo:traulls} is executed. Then 
	\begin{equation*}
		\lim\limits_{k\to \infty} \frac{\|\lambda_k\|}{\mu_k} = 0.
	\end{equation*}
\end{lemma}

In the spirit of~\cite[Lemma 4.3]{conn-etal:1996b}, we state the main convergence results.
\begin{lemma}\label{lemma:key_lemma}
	Suppose that assumptions~\ref{assumption:functions_C2},~\ref{assumption:feasible_linear_cons} hold. Let $\left(x_k\right)_{k\in \calK}$ be a sequence of iterates in $\Omega$ satisfying Assumption~\ref{assumption:iterates_domain_bounded} and converging to $x_*$ for which Assumption~\ref{assumption:cq_null_space_jacobian} holds. Let $\lambda_* = \lambda(x_*)$, $(\lambda_k)_{k \in \calK}$ be any sequence of vectors and $(\mu_k)_{k \in \calK}$ be a non-decreasing sequence of positive scalars. Assume that~\eqref{eq:algo_traulls_criticality_test} holds with $\lim\limits_{k \in \calK}\omega_k = 0$.
	
	Then, there are positive constants $\kappa_2, \kappa_3$ such that
	\begin{equation}\label{eq:bound_1st_order_mult}
		\|\bar{\lambda}_k - \lambda_* \| \le \kappa_2 \omega_k + \kappa_3 \|x_k-x_*\|,
	\end{equation}
	\begin{equation}\label{eq:bound_ls_mult}
		\|\lambda(x_k) - \lambda_* \| \le \kappa_3 \|x_k-x_*\|,
	\end{equation}
	and
	\begin{equation}\label{eq:bound_norm_cons}
		\|c(x_k)\| \le \kappa_2 \frac{\omega_k}{\mu_k} + \kappa_3 \frac{\|x_k-x_*\|}{\mu_k} + \frac{\|\lambda_k - \lambda_* \|}{\mu_k}.
	\end{equation}
	If, in addition, $c(x^*) = 0$, then $x_*$ is a first-order critical point for problem~\eqref{pb:cnls} with corresponding Lagrange multipliers $\lambda_*$ and with $\lim\limits_{k \in \calK} \bar{\lambda}_k = \lim\limits_{k \in \calK} \lambda(x_k) = \lambda_*$.
\end{lemma}

\begin{proof}
	By assumptions~\ref{assumption:functions_C2} and~\ref{assumption:cq_null_space_jacobian}, for $k$ sufficiently large, the matrix $\left(C_kN_*\right)^\dagger$ is well defined, bounded and converges to $\left(C_*N_*\right)^\dagger$. Therefore, there is a constant $\kappa_2$ such that
	\begin{equation*}
		\|(\left(C_kN_*\right)^\dagger)^\top \| \le \kappa_2.
	\end{equation*}
	Furthermore, by Lemma~\ref{lemma:bound_proj_algrad_sol_nullspace}, we have
	\[ \|N_*^\top \nabla_x\Phi_k\| = \left\| \proj{T(x_*)}{\nabla_x \Phi_k} \right\| \le  \omega_k,\]
	for $k \in \calK$ large enough. Thus we may deduce
	\begin{equation}\label{eq:bound_mult_1st_ls}
		\begin{split}
			\|\bar \lambda_k - \lambda(x_k) \| & = \|(\left(C_kN_*\right)^\dagger)^\top N_*^\top\nabla f_k + \bar\lambda_k\| \\
			& = \|(\left(C_kN_*\right)^\dagger)^\top (N_*^\top\nabla f_k + (C_kN_*)^\top\bar\lambda_k)\| \\
			& = \|((C_kN_*)^\dagger)^\top (N_*^\top  \nabla_x \Phi_k)\| \\
			& \le 	\|(\left(C(x)N_*\right)^\dagger)^\top \| \ \|N_*^\top \nabla_x\Phi_k\| \\
			& \le \kappa_2 \omega_k.
		\end{split}
	\end{equation}
	Furthermore, by the mean value theorem for vector-valued functions, we have
	\[\lambda(x_k) - \lambda_* = \int_{0}^{1} \nabla \lambda(x(t))(x_k-x_*)dt,\]
	where $x(t) = x_k + t(x_*-x_k)$. The integrand is bounded for $x_k$ sufficiently close to $x_*$, hence
	\[\| \lambda(x_k) - \lambda_*\| \le \kappa_3 \|x_k-x_*\|,\]
	for $k\in \calK$ sufficiently large and for some positive constant $\kappa_3$, giving~\eqref{eq:bound_ls_mult}.
	Then, combining the latter inequality and~\eqref{eq:bound_mult_1st_ls}, we may deduce that
	\begin{equation*}
		\begin{split}
			\|\bar \lambda_k - \lambda_* \| & \le \|\bar \lambda_k - \lambda(x_k) \| + \|\lambda(x_k) - \lambda_* \| \\
			& \le \kappa_2 \omega_k + \kappa_3 \|x_k-x_*\|,
		\end{split}
	\end{equation*}
	which is inequality~\eqref{eq:bound_1st_order_mult}. 
	Next, by expression~\eqref{eq:1st_order_mult}, we have $c(x_k) = (\bar \lambda_k - \lambda_k) / \mu_k$, which leads to
	\begin{equation*}
		\begin{split}
			\|c(x_k) \| & \le \frac{\|\bar \lambda_k - \lambda_* \|}{\mu_k} + \frac{\|\lambda_k - \lambda_* \|}{\mu_k} \\
			& \le \kappa_2 \frac{\omega_k}{\mu_k} + \kappa_3 \frac{\|x_k-x_*\|}{\mu_k} + \frac{\|\lambda_k - \lambda_* \|}{\mu_k},
		\end{split}
	\end{equation*}
	where we have used~\eqref{eq:bound_1st_order_mult}.
This concludes the first part of the lemma. 
	
	Since $(\omega_k)_k$ tends to $0$ and $(x_k)_{k\in \calK}$ converges to $x_*$ by assumption, inequalities~\eqref{eq:bound_1st_order_mult} and~\eqref{eq:bound_ls_mult} imply that both sequences $(\bar \lambda_k)_k$ and $(\lambda(x_k))_k$ converge to $\lambda_*$ for $k \in \calK$. Therefore, from identity~\eqref{eq:identity_al_lag_grad}, we obtain that the sequence of gradients $\nabla_x \Phi_k$ converge to $\nabla_x \calL(x_*,\lambda_*)$. 
	
	Finally, assume that $c(x^*) = 0$. Because~\eqref{eq:algo_traulls_criticality_test} holds, we have on the one hand that 
	\begin{equation}\label{eq:norm_proj_grad_to_0}
		\lim\limits_{k \in \calK} \|\proj{\Omega}{x_k-\nabla_x \Phi_k} - x_k\|= 0.
	\end{equation}
	On the other hand, the set $\Omega$ being closed and convex, the projection mapping $\proj{\Omega}{\cdot}$ is continuous, so $\proj{\Omega}{x_k-\nabla_x \Phi_k} \to \proj{\Omega}{x_*-\nabla_x \calL(x_*,\lambda_*)}$ as $k \in \calK$ increases. With~\eqref{eq:norm_proj_grad_to_0}, we obtain
	\[\proj{\Omega}{x_*-\nabla_x \calL(x_*,\lambda_*)} = x_*.\]
	Therefore, we have established that $(x_*,\lambda_*)$ is a KKT point and the lemma is proved.
\end{proof}

We can now establish the global convergence of Algorithm~\ref{algo:traulls}.
\begin{theorem}\label{theo:global_convergence_traulls}
	Suppose that assumptions~\ref{assumption:functions_C2},~\ref{assumption:feasible_linear_cons} hold. Let $\left(x_k\right)_{k}$ be a sequence of iterates generated by Algorithm~\ref{algo:traulls} satisfying Assumption~\ref{assumption:iterates_domain_bounded}. Further assume that there is a subsequence $(x_k)_{k \in \calK}$ converging to $x_*$, for which Assumption~\ref{assumption:cq_null_space_jacobian} holds and let $\lambda_* = \lambda(x_*)$.
	Then, conclusions of Lemma~\ref{lemma:key_lemma} hold.
\end{theorem}
\begin{proof}
	As in the proof of Lemma~\eqref{lemma:key_lemma}, the inequalities~\eqref{eq:bound_1st_order_mult},~\eqref{eq:bound_ls_mult},~\eqref{eq:bound_norm_cons} can be derived from the assumptions. All that remains is to show that $c(x^*) = 0$. There are two cases to distinguish. 
	
	First, if the sequence of penalty parameters is bounded above, then there is $k$ sufficiently large such that $\|c(x_k)\| \le \eta_k$ for all $k$. Since $\eta_k \to 0$, we get that $c(x_*) = 0$.
	
	Next, if the sequence of penalty parameters is unbounded, then by the update rule, we have that $(\mu_k)_k$ grows to $\infty$. Hence, by Lemma~\ref{lemma:behavior_multipliers}, the ratio $\lambda_k /\mu_k$ converges to zero. Therefore, by inequality~\eqref{eq:bound_norm_cons}, we also obtain that $c(x_*) = 0$. 
	In both cases, $x_*$ is feasible with respect to the nonlinear constraints and the remaining conclusions of Lemma~\ref{lemma:key_lemma} hold.
\end{proof}
Our convergence analysis is based on the criticality measure~\eqref{eq:algo_traulls_criticality_test} but, as we already mentioned in Section~\ref{sec:algorithm}, the latter is not available in closed form for a general polyhedron of the form~\eqref{eq:linear_constraints}. This is what justified our choice to monitor criticality in our algorithm through the reduced gradient norm $\|P_{T(x_k)}[\nabla_x\Phi_k]\|$, which only involves the projection onto the tangent space of the currently active bounds. The two measures are linked through the tangent cone $\calT_\Omega(x_k)$, which is the dual cone of $\calN_\Omega(x_k)$. Indeed, one always has \[\|\proj{\Omega}{x_k-\nabla_x\Phi_k}-x_k\|\le\|P_{\calT_\Omega(x_k)}[-\nabla_x\Phi_k]\|\] 
and whenever the multipliers associated with the bounds active at $x_k$ are of the correct sign the projection onto the cone reduces to the projection onto $T(x_k)$, so that (see~\cite[Section 7.2]{conn-etal:1996b})
\[\|P_{\calT_\Omega(x_k)}[-\nabla_x\Phi_k]\|=\|P_{T(x_k)}[\nabla_x\Phi_k]\|.\] 
Consequently, under the assumption that these sign conditions hold at every outer iterate, the reduced-gradient rule~\eqref{eq:criticality_cond_reduced_grad} certifies the criticality test~\eqref{eq:algo_traulls_criticality_test}, and Theorem~\ref{theo:global_convergence_traulls} applies to the iterates the implementation produces. Without this assumption the reduced gradient controls only the tangential component of $\nabla_x\Phi_*$, and the limit point may not be critical.